
\paragraph{\textbf{Corollary 5.4.}}
\SGARefTarget{sga3:xiv:5.4}
{\itshape Let \(G\) be a smooth connected algebraic group over an
algebraically closed field \(k\). Let \(H\) be a connected
algebraic subgroup such that \(\mathfrak h\) contains a Cartan
subalgebra of \(\mathfrak g\), or equivalently such that \(H\)
contains a subgroup of type \((\mathrm C)\) of \(G\). Then the
number of conjugates of \(H\) containing a regular element of
\(G(k)\) equals the number of conjugates of \(\mathfrak h\)
containing a regular element of \(\mathfrak g\).}

Indeed, let \(g\) be a regular element of \(G(k)\), and let \(C\)
be the unique Cartan subgroup of \(G\) containing \(g\)
\textup{(XIII~2)}. The conjugates of \(H\) that contain \(g\) are
exactly those that contain \(C\)
\textup{(\SGARefLink{sga3:xiii:corollary:2:8}{XIII~2.8}\textup{(b)})}.

Similarly, let \(x\) be a regular element of \(\mathfrak g\). If a
conjugate \(\mathfrak h'\) of \(\mathfrak h\) contains \(x\), then
\[
  \operatorname{ad}(x)_{\mathfrak g/\mathfrak h'}
\]
is injective by \SGARefLink{sga3:xiii:corollary:5:4}{XIII~5.4}, and therefore
\[
  \mathfrak h'\supset\Nil(x,\mathfrak g)=\mathfrak d.
\]
Thus the number of conjugates of \(\mathfrak h\) containing \(x\)
equals the number containing the Cartan subalgebra
\(\mathfrak d\).

The subalgebra \(\mathfrak d\) is the Lie algebra of a subgroup
\(D\) of type \((\mathrm C)\) of \(G\). We may suppose that
\(C\subset D\). It remains to show that
\[
  H\supset C\quad\Longleftrightarrow\quad
  \mathfrak h\supset\mathfrak d.
\]
By \SGARefLink{sga3:xiii:corollary:5:5}{XIII~5.5}, the inclusion
\(\mathfrak h\supset\mathfrak d\) implies \(H\supset D\), and
\emph{a fortiori} \(H\supset C\). Conversely, by \SGARefLink{sga3:xiv:5.1}{5.1} the Lie
algebra \(\mathfrak t\) of \(C\) contains a regular element \(x\)
of \(\mathfrak g\). Hence \(H\supset C\) implies
\(x\in\mathfrak h\), and, as observed above, this implies
\(\mathfrak h\supset\mathfrak d\). This completes the proof.

Let \(G\) be as in \SGARefLink{sga3:xiv:5.2}{5.2}, and suppose that the infinitesimal rank of
its fibers is locally constant, that is, condition
\((\mathrm C_0)\) holds. By \SGARefLink{sga3:xiv:5.2}{5.2}\textup{(c)}, there is a morphism of
functors on \((\operatorname{Sch}/S)^\circ\),
\[
  \mathcal C\longrightarrow\mathcal D,
\]
where
\[
  \begin{aligned}
  \mathcal C(S')&=
    \{\text{Cartan subgroups of }G_{S'}\},\\
  \mathcal D(S')&=
    \{\text{subgroups of type \((\mathrm C)\) of }G_{S'}\}.
  \end{aligned}
\]
By \SGARefLink{sga3:xiv:3.10}{3.10} and \SGARefLink{sga3:xiv:5.2}{5.2}\textup{(a)}, the functor \(\mathcal D\) is
represented by a smooth quasi-projective prescheme over \(S\).
Let \(D\) be the universal subgroup of type \((\mathrm C)\) of
\(G_{\mathcal D}\). One may then define the functor
\[
  \mathcal C_D\colon
  (\operatorname{Sch}/\mathcal D)^\circ
  \longrightarrow\operatorname{Ens}
\]
for the \(\mathcal D\)-group \(D\), exactly as \(\mathcal C\) was
defined for the \(S\)-group \(G\).

\paragraph{\textbf{Proposition 5.5.}}
\SGARefTarget{sga3:xiv:5.5}
{\itshape Under the preceding hypotheses, regard \(\mathcal C\) as
a functor over \(\mathcal D\). Then \(\mathcal C\) is
\(\mathcal D\)-isomorphic to the functor \(\mathcal C_D\) of
Cartan subgroups of \(D\).}

This follows immediately from the next corollary.

\paragraph{\textbf{Corollary 5.6.}}
\SGARefTarget{sga3:xiv:5.6}
{\itshape Let \(G\) be a smooth \(S\)-prescheme in groups of finite
presentation, and let \(D\) be a subgroup of type \((\mathrm C)\)
of \(G\). There is a bijective correspondence between the Cartan
subgroups of \(G\) contained in \(D\) and the Cartan subgroups of
\(D\). More precisely, for a subgroup \(H\subset D\), \(H\) is a
Cartan subgroup of \(G\) if and only if it is a Cartan subgroup of
\(D\).}

This is a special case of \SGARefLink{sga3:xii:7.9}{XII~7.9}\textup{\SGARefLink{sga3:xii:7.9:item:c:02}{(c)}}, taking into account
that, over an algebraically closed field, every subgroup of type
\((\mathrm C)\) of \(G\) contains a Cartan subgroup of \(G\).

For the next section, the principal result obtained here is
\SGARefLink{sga3:xiv:5.2}{5.2}\textup{(c)} with \(H\) a Cartan subgroup of \(G\). It permits
the formulation of \SGARefLink{sga3:xiv:5.5}{5.5} and thus gives a useful dévissage of
\(\mathcal C\).
